% Source: https://artofproblemsolving.com/community/c3h1227228_are_they_correct
% Stable identity: aops-topic-1227228-q02
The $13$ths of the twelve months fall on weekdays shifted from the $13$th of January by the cumulative month lengths modulo $7$. For a common year the shifts are
\[0,3,3,6,1,4,6,2,5,0,3,5\]
(Jan, \dots, Dec); for a leap year they are $0,3,4,0,2,5,0,3,6,1,4,6$. In the common year the residues $0,3,3,6,1,4,6,2,5,0,3,5$ occur with multiplicities: $3$ appears three times, every residue at least once. In the leap year: $0$ three times, every residue at least once. So in every year each weekday is the $13$th of at least one month, and the maximum is $3$. Hence the largest number of Friday the $13$ths is $\boxed{3}$ and the smallest is $\boxed{1}$.
